% !Mode:: "TeX:UTF-8"
% +-----------------------------------------------------------------------------
% | File: resume-zh.tex
% | Author: huxuan
% | E-mail: i(at)huxuan.org
% | Created: 2012-12-18
% | Last modified: 2020-04-20
% | Description:
% |     A Chinese Resume Example in LaTeX based on resumecls
% |
% | Copyright (c) 2012-2020 by huxuan. All rights reserved.
% +-----------------------------------------------------------------------------

\documentclass[color]{resumecls}

\include{resume-common}

\name{陈聪涌}
\organization{上海交通大学}
\address{上海市闵行区东川路800号，200240}
\mobile{+86 156 1177 5012}
\mail{619455181@qq.com}
\homepage{http://github.com/ndsf}
\leftfooter{最后更新: \today}
\rightfooter{\url{http://github.com/ndsf}}

\ctexset{today=small}

\begin{document}
\begin{table}
\maketitle
%%%%%%%%%%%%%%%%%%%%%%%%%%%%%%%%%%%%%%%%%%%%%%%%%%%%%%%%%%%%%%%%%%%%%%%%%%%%%%%
%\heading{经历(紧凑版)}
%\entry{0em}{Xlr}{{\bfseries 组织} \quad 部门 & 职位 & 起止时间}
%\entry{2em}{X}{%
    %描述1 \\
    %描述2 \\
%}
%%%%%%%%%%%%%%%%%%%%%%%%%%%%%%%%%%%%%%%%%%%%%%%%%%%%%%%%%%%%%%%%%%%%%%%%%%%%%%%
\heading{教育经历}
\entry{0em}{Xrr}{{\bfseries 上海交通大学} & 上海 & 2021.9 - 预计2026.6}
\entry{2em}{lXX}{%
    博士研究生在读 \quad 计算机学院 \quad 计算机科学与技术专业 \quad 内存存储系统、分布式系统方向 \\
    % 技能 \quad C, C++, RDMA, CXL, DPU, Linux内核, 文件系统, 索引, 持久内存 \\
    
    %学位2 & 学院2 & 专业2 \\
}
\entry{0em}{Xrr}{{\bfseries 北京航空航天大学} & 北京 & 2017.9 - 2021.9}
\entry{2em}{lXX}{%
    学士学位 & 软件学院 & 软件工程专业 \\
    
    %学位2 & 学院2 & 专业2 \\
}
% \entry{2em}{lXX}{%
%     主要课程 &  	软件工程基础(95)\ 	数据管理技术(96)\ 面向对象程序设计(JAVA)(94)\ 程序设计基础训练(100) \ {离散数学(信息类)(96)}\  	数据结构与程序设计(信息类)(99)\ 面向程序设计的硬件基础(96)\ 	离散数学（2）(96)\  	算法分析与设计 (100)\ 	工科数学分析(1)(94)\ 	工科数学分析(2)(89)\  {概率统计A(97)}\ 计算机网络 (94)\\
% }
\entry{2em}{lXX}{%
     GPA &  	平均GPA为3.6835，年级排名为17/167，只计算专业课程排名为5/167。% 2017~2018学年GPA为3.79/4.00，年级排名为5/164，2018~2019学年GPA为3.78，年级排名为12/167，2019~2020上半学年GPA为3.88/4.00。由于只统计专业课程成绩，学年GPA高于平均GPA。\\
}
\entry{2em}{lXX}{%
     外语 & 四级成绩为596分，六级成绩为616分。 \\
}
%%%%%%%%%%%%%%%%%%%%%%%%%%%%%%%%%%%%%%%%%%%%%%%%%%%%%%%%%%%%%%%%%%%%%%%%%%%%%%%
\heading{科研经历}
\entry{0em}{Xr}{{\bfseries 上海交通大学} & 2021.9 - 至今}
\entry{2em}{X}{数据驱动软件技术实验室 \quad 博士研究生 \quad 内存存储系统、分布式系统方向 \quad 导师黄林鹏教授}
\entry{2em}{X}{% 分离式内存架构将计算和内存资源分成独立的池，通过RDMA、CXL等高速网络连接，可以提高资源利用率。
    参与项目“\textbf{分离式内存（DM）计算框架研究}”。面向RDMA、CXL等高速网络设计高并发分布式哈希索引。现有索引受限于网络原语的有限功能，在写入期间存在写入位置查找、并发冲突、调整大小等待、缓存一致性等开销，降低高并发和偏斜场景性能。本项目通过重新设计索引结构、利用共享接收队列等RDMA特性、计数器触发合并、乐观调整大小检测等方式，基于RDMA、C++实现高性能哈希索引。以第1作者身份正在投稿论文。针对现有DM范围索引依赖节点锁导致的写入位置查找及热点页写竞争，以及传统无锁Bw-tree在DM上存在指针追踪等问题，设计面向DM的Bw-tree。通过顺序追加式缓冲区、多区域指纹过滤、就地增量合并以及页面读写聚合，基于RDMA/CXL实现无锁范围索引。以第 1 作者身份正在投稿论文。\\
    % 《Redesigning Hash Indexes on Disaggregated Memory for Minimal Network Overhead》\cite{chen2026rdma}到ICDE 2026（CCF-A会议）。 \\
}
\entry{2em}{X}{%
    参与CCF-华为胡杨林基金系统软件专项项目“\textbf{面向新介质的高性能高可靠文件系统}"。面向持久CPU缓存环境（如Intel eADR、CXL GPF），基于openEuler的EulerFS研究Linux内核文件系统。以第1作者身份发表文章《Redesigning Data and Metadata Updates in PM File Systems with Persistent CPU Caches》 \cite{chen2024redesigning}于DASFAA 2024（CCF-B会议）。本文提出完全去除缓存行刷新会因随机缓存行驱逐导致写放大，并通过监控系统调用和内存映射的文件访问模式，选择最有效的更新机制。以第1作者身份发表文章《FusionFS: A Contention-Resilient File System for Persistent CPU Caches》 \cite{chen2025fusionfs}于ACM TACO 2025（CCF-A期刊）。本文进一步分类并使用缓存分配等技术缓解持久CPU缓存面临的多种争用，基于C语言实现Linux内核文件系统FusionFS \url{https://github.com/SJTU-DDST/FusionFS}。发明专利《索引方法、程序产品、介质以及电子设备》（CN119493769A）。 \\
    %这里可以使用bibtex，如\cite{label} \\
}
%%%%%%%%%%%%%%%%%%%%%%%%%%%%%%%%%%%%%%%%%%%%%%%%%%%%%%%%%%%%%%%%%%%%%%%%%%%%%%%
% \heading{工作经历}
% \entry{0em}{Xr}{{\bfseries 旷视科技} & 2019.7 - 2019.8}
% \entry{2em}{X}{IC Team \quad 算法实习生}
% \entry{4em}{X}{%
%     参与数据获取、量化模型训练等工作。 \\
% }
% %%%%%%%%%%%%%%%%%%%%%%%%%%%%%%%%%%%%%%%%%%%%%%%%%%%%%%%%%%%%%%%%%%%%%%%%%%%%%%%
% \heading{获得荣誉}
% \entry{0em}{Xr}{%
%     北京航空航天大学2019至2020学年学习优秀奖学金一等奖 & 2020.12 \\
% 	北京航空航天大学2018至2019学年学习优秀奖学金一等奖 & 2019.12 \\
% 	北京航空航天大学2017至2018学年学习优秀奖学金一等奖 & 2018.12 \\
% 	第十二届北航程序设计竞赛三等奖 & 2016.12 \\
% 	第二十一届全国青少年信息学奥林匹克联赛(NOIP)省一等奖 & 2015.12 \\
% }
%%%%%%%%%%%%%%%%%%%%%%%%%%%%%%%%%%%%%%%%%%%%%%%%%%%%%%%%%%%%%%%%%%%%%%%%%%%%%%%
% \heading{专业技能}
% \entry{0em}{lX}{%
%     熟悉 & Javascript \\
%     掌握 & C/C++, Java, Python, Git \\
%     使用 & \LaTeX \\
% }
%%%%%%%%%%%%%%%%%%%%%%%%%%%%%%%%%%%%%%%%%%%%%%%%%%%%%%%%%%%%%%%%%%%%%%%%%%%%%%%
%\heading{其他列举事项-如个人爱好，网络资料等}
%\entry{0em}{lX}{%
    %标签1 & 标签对应内容 \\
    %标签2 & 标签对应内容 \\
%}
%%%%%%%%%%%%%%%%%%%%%%%%%%%%%%%%%%%%%%%%%%%%%%%%%%%%%%%%%%%%%%%%%%%%%%%%%%%%%%%
% 如果不需要发表成果，注释这一段即可。
% \heading{发表成果}
% \vspace{-6em}
% \bibliography{resume}
%%%%%%%%%%%%%%%%%%%%%%%%%%%%%%%%%%%%%%%%%%%%%%%%%%%%%%%%%%%%%%%%%%%%%%%%%%%%%%%
% 两页或者多页的示例。
\end{table}
\clearpage
\begin{table}

\heading{科研经历}
\entry{0em}{Xr}{{\bfseries 北京航空航天大学} & 2019.9 - 2019.11}
%\entry{2em}{X}{实验室 \quad 职位}
\entry{2em}{X}{%
    参与项目“压缩视频质量增强”，具体负责“复现部分对比实验代码、比较结果”。以第4作者身份参与文章《Multi-level Wavelet-based Generative Adversarial Network for Perceptual Quality Enhancement of Compressed Video》 \cite{wang2020multi}，发表于ECCV 2020（CCF-B会议）。 \\
    %这里可以使用bibtex，如\cite{label} \\
}

\heading{工作经历}
\entry{0em}{Xr}{{\bfseries 旷视科技} & 2019.7 - 2019.8}
\entry{2em}{X}{IC Team \quad 算法实习生}
\entry{4em}{X}{%
    参与数据获取、量化模型训练等工作。 \\
}
%%%%%%%%%%%%%%%%%%%%%%%%%%%%%%%%%%%%%%%%%%%%%%%%%%%%%%%%%%%%%%%%%%%%%%%%%%%%%%%
\heading{获得荣誉}
\entry{0em}{Xr}{%
    北京航空航天大学2019至2020学年学习优秀奖学金一等奖 & 2020.12 \\
	北京航空航天大学2018至2019学年学习优秀奖学金一等奖 & 2019.12 \\
	北京航空航天大学2017至2018学年学习优秀奖学金一等奖 & 2018.12 \\
	第十二届北航程序设计竞赛三等奖 & 2016.12 \\
	第二十一届全国青少年信息学奥林匹克联赛(NOIP)省一等奖 & 2015.12 \\
}

\heading{发表成果}
\vspace{-6em}
\bibliography{resume}

%\heading{第二页小标题}
%\entry{0em}{X}{
    %每个页面都包含在一个 table 环境中。 \\
    %表格之间的 \verb|\clearpage| 会新建一个页面。 \\
%}
%%%%%%%%%%%%%%%%%%%%%%%%%%%%%%%%%%%%%%%%%%%%%%%%%%%%%%%%%%%%%%%%%%%%%%%%%%%%%%%
% \heading{校园经历}
% \entry{0em}{Xr}{%
% 	1. 基于Java的跨平台小游戏。为程序设计实践项目，使用Java语言基于libGDX开发，实现游戏功能和跨桌面端和移动端的跨平台功能。独自开发，负责软件设计、软件开发、软件测试、文档撰写等工作。 & 2018.9\\
% 	2. 小型二手货交易系统。为软件工程课程项目，使用Python语言基于Django开发，实现账号系统、商品系统、后台管理等功能。在项目中作为负责人和三名组员共同开发，负责需求调研、软件设计、软件开发、软件测试、文档撰写等工作。 & 2019.4 - 2019.6\\
% 	3. 校园网客户端。为面向对象程序设计(JAVA)课程项目，使用Java语言基于JavaFX开发，实现GUI、保持连接、数据统计等功能。在项目中作为负责人和一名组员共同开发，负责软件设计、软件开发、软件测试、文档撰写等工作。 & 2019.4 - 2019.6\\
% 	4. 书籍影视交流平台。为软件工程实践（一）课程项目，使用Javascript语言基于Express、GraphQL、React、MongoDB进行敏捷开发，实现帐号系统、书籍交流、影视交流、小组系统、后台管理等功能。在项目中作为负责人和三名组员共同开发，负责需求调研、软件设计、软件开发、软件测试、文档撰写等工作。 & 2019.8\\
% 	5. 科技专家资源共享平台。为软件系统分析与设计课程项目，使用Javascript语言基于Express、GraphQL、Vue.js、MongoDB、ElasticSearch进行开发，实现帐号系统、资源爬虫、专家主页、主页认领、资源检索、成果转让、资源统计、资源变现等功能。在项目中作为组员和其余八名组员共同开发，负责需求模型、设计数据库模型、后端的架构设计和开发。 & 2019.9 - 2019.12\\
% 	6. 基于P2P的Android聊天应用。为软件设计模式课程项目，使用Javascript和Java语言基于Express、GraphQL、Android进行开发。在项目中作为组员和其余八名组员共同开发，负责GraphQL相关代码编写。 & 2019.9 - 2019.12\\
% 	7. C0编译器。为编译原理课程项目，使用C++进行开发。 & 2019.12\\
% 	8. 在北京航空航天大学李红裔教授的人工智能讨论班中学习。 & 2019.3 - 至今\\
% }
\end{table}
\end{document}
